%=====================================================================
\chapter{The Boundary With Other Courses}\label{app:boundary}
%=====================================================================
\normalfigures

This course sits downstream of four others and alongside a fifth. It assumes a
good deal and repeats a little, and both the assumptions and the repetitions
are deliberate. This appendix records what belongs where, so that a lecturer
can see what to assume and a student can see why a familiar topic reappears.

\section{The Edition History, Which Is Unusual}

Before the boundaries, the provenance --- because this course's position in the
degree is less settled than any of its siblings'.

\rowcolors{2}{white}{lightbg}
\begin{longtable}{@{}L{1.8cm}L{5.0cm}L{3.0cm}L{4.6cm}@{}}
\toprule
\rowcolor{primary!15}
\textbf{Edition} & \textbf{Title} & \textbf{Where} & \textbf{Content
published?} \\
\midrule
\endhead
\textbf{2017} & E-Commerce & BSSE elective, p.~38 & \textbf{Yes --- p.~94, the
only full specification in any edition} \\
2023 & --- & \textbf{Absent entirely} & One incidental mention, p.~23 \\
2025 & Digital Marketing \& E-Commerce & Interdisciplinary / allied list,
item 11 & No --- title and credit hours only \\
\bottomrule
\end{longtable}

\smallskip
So the single published outline for this course is a \textbf{2017 specification
for an elective in a different degree}, the current editions do not specify it
at all, and the department teaches it as \textbf{INFT-4041, BSIT, 5th
semester, 3 credit hours}. Everything in this handout is traced to that 2017
outline, with each addition justified in the Course Specification at the front.

\begin{alertbox}[A prerequisite under another name]
HEC's 2017 entry gives the prerequisite as \textbf{Web Engineering}. The
department has no course of that name; its equivalent is \textbf{CS-207 /
INFT-4005 Web Technologies} at 4th semester.

\smallskip
\textbf{The name differs and the dependency is satisfied.} Recorded here rather
than silently matched, because a reviewer comparing the two documents will
notice, and the answer should be written down before they ask.
\end{alertbox}

\section{Web Technologies --- the Prerequisite}

\begin{conceptbox}[The rule]
\textbf{Web Technologies owns the request. This course owns what is at the
other end of it.}

\smallskip
A student arriving here should be able to read an HTTP request and response,
write HTML and CSS, and know what a form submission is. What this course adds
is \emph{an application that answers}, with a database behind it, money passing
through it, and consequences attached to getting it wrong.
\end{conceptbox}

\rowcolors{2}{white}{lightbg}
\begin{longtable}{@{}L{6.8cm}L{3.4cm}L{4.2cm}@{}}
\toprule
\rowcolor{primary!15}
\textbf{Topic} & \textbf{Owned by} & \textbf{Here} \\
\midrule
\endhead
HTTP methods, status codes, headers & Web Technologies & Assumed; used in
\chref{security} and \chref{seo} \\
HTML, CSS, responsive layout & Web Technologies & Assumed; \chref{ux} is about
\emph{overriding} templates, not writing CSS \\
Cookies, as a browser feature & Web Technologies & \chref{basket} asks where
the state they point at actually lives \\
Forms and validation & Web Technologies & \chref{checkout} is about the
\emph{sequence} of forms and what each commits \\
A server-side language & Web Technologies & Appendix~B is a refresher, not a
first course \\
\bottomrule
\end{longtable}

\section{Database Systems}

\begin{conceptbox}[The rule]
\textbf{Database Systems owns the model. This course owns the consequences of
somebody else's model.}

\smallskip
You will not design a schema here. You will read one you did not write, find
out what it costs, and discover that the answer changes with its size.
\end{conceptbox}

\begin{tight}
  \item \textbf{Assumed:} SQL selects and joins, primary and foreign keys,
        normalisation, what an index is.
  \item \textbf{Added here:} reading the query a screen \emph{actually} ran
        (\chref{catalogue}); measuring an index rather than assuming it
        (\chref{catalogue} found two that made the page slower at 50{,}000
        products); and a schema whose shape is a historical record of
        somebody's decisions --- 18 product tables, options that grow by
        addition and variants that grow by multiplication
        (\chref{variations}).
  \item \textbf{The useful repetition:} every Database Systems student has been
        told that indexes make queries faster. \chref{catalogue} is where that
        becomes a measurement with a crossover point in it.
\end{tight}

\section{Information Security}

\begin{conceptbox}[The rule]
\textbf{Information Security owns the principles. \chref{security} owns one
shop, measured twice.}

\smallskip
That chapter is not a survey of attacks. It is a battery of fifteen checks run
against a real installation before and after one configuration file, and its
central finding --- that the famous attacks were already handled and every real
problem was in the deployment --- is one a principles course cannot deliver.
\end{conceptbox}

\begin{tight}
  \item \textbf{Assumed:} the CIA triad, authentication versus authorisation,
        hashing versus encryption, what injection and cross-site scripting
        \emph{are}.
  \item \textbf{Added here:} measuring them (\chref{security}); the cost factor
        as a trade-off with a number attached (\chref{accounts}); a timing side
        channel caused by the defence itself; and \textbf{the legal boundary},
        which is stated in Appendix~A, in \chref{security} and again in
        \chref{legal}.
  \item \textbf{Not covered here:} cryptographic construction, network
        security, forensics. Those stay with Information Security.
\end{tight}

\section{Software Engineering}

\begin{tight}
  \item \textbf{Assumed:} requirements, version control, testing, the idea that
        a system has a lifecycle.
  \item \textbf{Added here:} the build-versus-buy-versus-extend decision made
        against a real platform (\chref{planning}); extension as the discipline
        that keeps a system upgradable (\chref{modules}); and deployment
        measured rather than described (\chref{deploy}).
  \item \textbf{The overlap worth naming:} \chref{planning}'s decision is a
        software engineering decision, taught here because it is the one that
        determines everything after it in this course. A student who has made
        it in Software Engineering should make it again here with a specific
        platform in front of them, which is a different exercise.
\end{tight}

\section{Virtual Systems and Services --- the Same Semester}

\begin{alertbox}[This overlap needs stating, because the two courses are taught
side by side]
\textbf{INFT-4008 Virtual Systems and Services} is also a 5th-semester BSIT
course, and it also uses containers. A student taking both will meet Docker
twice in one semester.

\smallskip
\textbf{The split.} Virtual Systems and Services owns \emph{virtualisation as a
subject}: hypervisors, container internals, orchestration, cloud service
models, how a container differs from a virtual machine at the kernel level.

\smallskip
\textbf{This course uses containers as a tool and says so.} \chref{stack} spends
one chapter on what a container is --- enough to use Compose confidently and no
more --- because its purpose is a reproducible shop, not an education in
virtualisation. \chref{deploy} returns to deployment, again from the
application's side: saturation, backup, restore.

\smallskip
\textbf{For the lecturer:} if the two courses are taught in the same semester,
\chref{stack} can be shortened by agreement, and the time spent on
\chref{catalogue}, which always needs more. If they are not, \chref{stack}
stands alone as written.
\end{alertbox}

\section{What This Course Does Not Teach}

Stated so that nobody is surprised at the end of it.

\begin{tight}
  \item \textbf{Front-end frameworks.} No React, no Vue. \chref{ux} overrides
        templates in the platform's own engine.
  \item \textbf{Writing a cart from scratch.} \chref{planning} argues
        explicitly against it, and the argument is the lesson.
  \item \textbf{Digital marketing as a discipline.} \chref{seo} covers what is
        checkable and says plainly what it is leaving out, which matters because
        the 2025 edition renames the course.
  \item \textbf{Cloud deployment on a named provider.} \chref{deploy} measures
        a deployment; which cloud it sits on is a different course's subject
        and dates faster than this handout.
  \item \textbf{Payment integration with a live provider.} \chref{gateway}
        ships its own, because a merchant account cannot be a prerequisite for
        a laboratory.
\end{tight}
